\documentclass[10pt,a4paper]{amsart}
\usepackage[margin=19mm]{geometry}
\usepackage{amsmath,amssymb,mathtools}
\usepackage[scr=rsfs,cal=euler]{mathalfa}
\usepackage{xfrac}
\usepackage[unicode,hidelinks,bookmarks=false]{hyperref}
\newcommand{\ie}{i.e.\ }
\newcommand{\cf}{cf.\ }
\newcommand{\resp}{respectively\ }
\newcommand{\Subsection}{\subsection}
\providecommand{\nobreakdash}{\nobreak\hbox{-}}
\setlength{\emergencystretch}{2em}
\multlinegap0pt
\usepackage{xcolor}
\usepackage{amssymb}
\usepackage{tikz}
\usepackage{float}
\usepackage{algorithm}\usepackage{algpseudocode}
\newtheorem{lemma}{Lemma}
\title{Endnotes using pagenote package at the end of each chapter}
\author{Christine T. Cheng, Chelsea Ann Lambert}
\date{Source: arXiv:2601.00957v1 (CC BY 4.0)}
\begin{document}
\maketitle

\noindent\emph{Source and licence.} Endnotes using pagenote package at the end of each chapter. Version arXiv:2601.00957v1 (CC BY 4.0), distributed by arXiv under the Creative Commons Attribution 4.0 International licence, http://creativecommons.org/licenses/by/4.0/, as recorded in the arXiv licence metadata for this exact version and shown on the abstract page https://arxiv.org/abs/2601.00957v1. Full licence notice: \texttt{licenses/arxiv-2601.00957-CC-BY-4.0.txt}.\par\smallskip

\noindent\textit{Editorial scope.} Counted unit: the article's lemma lemmaWhitman2 (source0.tex lines 1084--1087). The article's complete original proof of it is delivered with it (source0.tex lines 1090--1103, one \texttt{proof} environment of 14 lines). No mathematics is written, repaired or re-derived: the statement and the proof are the article's own lines, in the article's own notation, and no retained line is substituted or re-flowed.  No further source-local context is needed: every cross-reference of every retained line resolves inside the two delivered spans.  Nothing was omitted from any retained span, so the package carries no omission record.  No retained line contains a citation, so the wrapper carries no bibliography and no bibliography entry.  No retained line contains a figure environment, an image-inclusion command or drawing code, so no image is delivered and the package declares no asset dependency.  Editorial formatting only, in addition: an AMS \texttt{amsart} preamble that restates the article's own declarations and macros used by the retained lines, namely the environments \texttt{lemma} on separate counters, together with the standard packages those lines need.  The retained environments print in this document as Lemma 1 in order of appearance, whereas the article prints the same items with the numbering of its own full document.  The complete native source of the article as distributed in the arXiv e-print for this exact version is bundled unchanged as \texttt{sources/192039/source0.tex}.
\begin{lemma}
Let $G= (G_k , A_k, B_k) \circ (G_{k-1}, A_{k-1}, B_{k-1})  \circ \hdots \circ (G_1, A_1, B_1) \circ G_0$. Then it is the case that $\alpha(G) = |B_k| + |B_{k-1}|+ \hdots + |B_1| + \alpha(G_0)$ and $\beta(G) = |A_k| + |A_{k-1}|+ \hdots + |A_1| + \beta(G_0)$.
\label{lemmaWhitman2}
\end{lemma}

\begin{proof}
We prove the formula for $\alpha(G)$ first using induction on the number of components of $G$.   When $G$ has only one component, $G = G_0$ so $\alpha(G) = \alpha(G_0)$.   Assume that the result holds for graphs that have $k$ or fewer components.  Let $G$ be a graph with $k+1$ components so $G= (G_k , A_k, B_k) \circ (G_{k-1}, A_{k-1}, B_{k-1})  \circ \hdots \circ (G_1, A_1, B_1) \circ G_0$. Let $H = (G_{k-1}, A_{k-1}, B_{k-1})  \circ \hdots \circ (G_1, A_1, B_1) \circ G_0$ so $G = (G_k, A_k, B_k) \circ H$.   Denote as $B_H$   a maximum-sized stable set in $H$.  Clearly, $B_k \cup B_H$ is a stable set of $G$ so $\alpha(G) \ge |B_k| + |B_H|$.  We will now show that $\alpha(G) \leq |B_k| + |B_H|$.

 Let $D$ be a stable set of $G$,  $D_k = D \cap V(G_k)$ and $D_H = D \cap V(H)$.  Notice that both $D_k$ and $D_H$ must be stable sets of $G_k$ and $H$ respectively.  Furthermore,  $D_k$ can contain at most one vertex in $A_k$ since the vertices in $A_k$ induce a clique.  Let us consider two cases:  (i) $|D_k \cap A_k| = 0$ and (ii) $|D_k \cap A_k| = 1$.  For case (i),  $|D| = |D_k| + |D_H| \leq |B_k| + |B_H|$.    For case (ii), since $D_k \cap A_k \neq \emptyset$,  $D_H$ has to be empty because  every vertex in $A_k$ is adjacent to every vertex of $H$ in $G$.  Thus,  $|D| = |D_k|  \leq |B_k| + 1  \leq |B_k| + |B_H|$.  The latter inequality holds because $H$ is a non-empty graph so $|B_H| \ge 1$.  We have now shown that for cases (i) and (ii), every stable set of $G$ has size at most $|B_k| + |B_H|$.  It follows that $\alpha(G) = |B_k| + |B_H|$.  But by the induction hypothesis, $|B_H| = |B_{k-1}|+ \hdots + |B_1| + \alpha(G_0)$ so $\alpha(G) = |B_k| + |B_{k-1}|+ \hdots + |B_1| + \alpha(G_0)$. 
 
 Now, $\alpha(G) + \beta(G) = |V(G)|$ and $|V(G)| = |V(G_k)| + |V(G_{k-1})| + \hdots +  |V(G_0)|$.   Thus, 
 \begin{eqnarray*}
 \beta(G) & = &  |V(G)| - \alpha(G) \\
                & = & (|V(G_k)| + |V(G_{k-1})| + \hdots + |V(G_1)| + |V(G_0)|) - (|B_k| + |B_{k-1}|+ \hdots + |B_1| + \alpha(G_0)) \\
                & = & (|V(G_k)| - |B_k|) + (|V(G_{k-1})| - |B_{k-1}|) +  \hdots + (|V(G_1)| - |B_1|) + (|V(G_0)| - \alpha(G_0)) \\
                & = &  |A_k| + |A_{k-1}| + \hdots + |A_1| + \beta(G_0)
 \end{eqnarray*}
 where the last equality uses the fact that $\alpha(G_0) + \beta(G_0) = |V(G_0)|$.
\end{proof}

\end{document}
